%% -*- coding: utf-8 -*-
%% Migrated from paper/manuscript/src/results.tex

\section{Results}

We developed a GPU-accelerated computational framework for phase-amplitude coupling (PAC) analysis implemented in PyTorch. The framework provides substantial speedups over existing CPU-based implementations while maintaining calculation accuracy, and introduces trainable PAC parameters for end-to-end optimization through gradient descent.

\subsection{PAC Value Confirmation with an Existing Package}
To validate the accuracy of our GPU-accelerated implementation, we compared PAC values computed by our method (gPAC) against TensorPAC \cite{combrisson_tensorpac_2020}, an established CPU-based PAC analysis package (Figure~\ref{fig:00_pac_value_comparision}). Using identical synthetic signals with known coupling parameters, we observed near-identical PAC values between both implementations. The mean squared error between gPAC and TensorPAC outputs was negligible, confirming that GPU acceleration does not compromise calculation accuracy. Additionally, the trainable version of gPAC produced equivalent PAC values prior to parameter optimization, demonstrating that the differentiable architecture preserves numerical fidelity.

\subsection{Speed Comparison with an Existing Package}
We systematically benchmarked computation speed across multiple parameter dimensions to characterize the performance advantage of GPU-accelerated PAC computation.

\paragraph{Batch size.} Increasing batch size from $2^3$ to $2^6$ demonstrated the parallel processing advantage of our GPU implementation (Figure~\ref{fig:01_batch_size}). While TensorPAC processing time scaled linearly with batch size on CPU, gPAC maintained near-constant computation time on GPU due to efficient parallelization of independent samples.

\paragraph{Chunk size.} The effect of temporal chunking on computation speed revealed that our implementation efficiently handles varying segment lengths (Figure~\ref{fig:02_chunk_size}), with GPU acceleration providing consistent speedups across all chunk sizes tested ($2^0$ to $2^3$).

\paragraph{Number of channels.} Scaling from $2^3$ to $2^6$ channels showed that gPAC achieves substantial speedups with increasing channel counts (Figure~\ref{fig:03_n_chs}), as multi-channel processing is inherently parallelizable on GPU architectures.

\paragraph{Signal duration.} Processing times for signal durations from $2^0$ to $2^3$ seconds demonstrated that GPU acceleration provides greater relative speedup for longer recordings (Figure~\ref{fig:04_t_sec}), which is particularly relevant for clinical and research applications involving continuous recordings.

\paragraph{Sampling frequency.} Benchmarking at $2^9$ Hz and $2^{10}$ Hz sampling rates showed that higher sampling frequencies, which produce larger data volumes, benefit proportionally more from GPU parallelization (Figure~\ref{fig:05_fs}).

\paragraph{Number of frequency bands.} The computational advantage of gPAC became particularly pronounced when computing PAC across many frequency band combinations (Figure~\ref{fig:07_pha_n_bands}). Testing from 10 to $10^2$ phase frequency bands revealed that the GPU implementation scales favorably with the number of frequency band pairs analyzed.

\paragraph{Permutation testing.} Surrogate-based statistical testing is computationally expensive in PAC analysis. Our implementation parallelizes permutation computations on GPU, achieving substantial speedups for $2^0$ to $2^2$ permutations (Figure~\ref{fig:08_n_perm}).

Overall, across all parameter combinations tested, gPAC achieved approximately 100-fold speedup compared to TensorPAC, with the largest gains observed in high-dimensional analyses involving many channels, frequency bands, and permutations.

\subsection{Trainable Phase-Amplitude Coupling}
A key innovation of our framework is making PAC parameters fully differentiable, enabling training through backpropagation algorithms. Specifically, we implemented: (i) learnable bandpass filter parameters for data-driven frequency band selection, (ii) a differentiable Hilbert transformation preserving gradient flow, and (iii) adaptable phase-amplitude binning for modulation index calculation. To demonstrate the utility of trainable PAC, we trained a model incorporating the trainable PAC module for a classification task distinguishing between coupled and uncoupled oscillatory signals. The framework achieved high classification accuracy on synthetic data and successfully identified physiologically relevant theta-gamma coupling patterns, suggesting that end-to-end optimization can discover meaningful coupling relationships from neural recordings.

%%%% EOF
